\documentclass[11pt]{article}
\usepackage[utf8]{inputenc}\usepackage[T1]{fontenc}\usepackage[french]{babel}
\usepackage{amsmath,amssymb,booktabs}\usepackage[margin=2.2cm]{geometry}
\begin{document}
\section*{Densité attendue de raies moléculaires pour $^{87}$Sr}

\paragraph{1. Quelle loi de LeRoy-Bernstein ?}
$E_\nu = \big[(\nu_D-\nu)\,H_n\big]^{2n/(n-2)}$, soit $E^{1/6}$ linéaire en $\nu$ pour $n=3$ et $E^{1/3}$ linéaire pour $n=6$.
Test sur les niveaux $0_u$ mesurés de $^{88}$Sr (Zelevinsky \emph{et al.}) :
\begin{center}
\begin{tabular}{lcccc}\toprule
$\nu$ & $-1$ & $-2$ & $-3$ & $-4$\\ \midrule
$E$ (MHz) & 0.435 & 23.932 & 222.161 & 1084.093\\
$E^{1/6}$ & 0.870 & 1.698 & 2.461 & 3.205 \\
$E^{1/3}$ & 0.758 & 2.882 & 6.057 & 10.273\\ \bottomrule
\end{tabular}
\end{center}
Les écarts successifs de $E^{1/6}$ sont presque constants ($0.83,\ 0.76,\ 0.74$), ceux de $E^{1/3}$ ne le sont pas ($2.1,\ 3.2,\ 4.2$).
Près du seuil, le potentiel excité suit donc la loi en $C_3/r^3$ ($n=3$), et non en $C_6/r^6$ :
\[
H_3(^{88}\text{Sr}) = 0.777~\text{MHz}^{1/6}.
\]
L'expression $3H_6\epsilon^{2/3}$ utilisée dans le texte est l'espacement $dE/d\nu$ d'un potentiel en $C_6$ : ni la bonne loi, ni la bonne grandeur.

\paragraph{2. Passage à $^{87}$Sr.} $C_3$ est électronique (identique pour les deux isotopes) et $H_3\propto\mu^{-1/2}$ :
\[
H_3(^{87}\text{Sr}) = H_3(^{88}\text{Sr})\sqrt{\frac{\mu_{88}}{\mu_{87}}} = 0.781~\text{MHz}^{1/6}.
\]

\paragraph{3. Nombre de niveaux vibrationnels par canal.}
Le nombre de niveaux liés à moins de $E$ sous le seuil d'un canal vaut environ $E^{1/6}/H_3$ :
\begin{center}
\begin{tabular}{lc}\toprule
Domaine sous le seuil & Niveaux par canal \\ \midrule
$0$--$266$~MHz & $\simeq 3.3$\\
$0$--$380$~MHz (jusqu'à la raie la plus profonde observée) & $\simeq 3.5$\\ \bottomrule
\end{tabular}
\end{center}
Ordre de grandeur des énergies de liaison successives (pour $\nu_D$ entier) : $0.2$, $15$, $170$, $930$~MHz. Il y a donc environ trois niveaux par canal dans la zone explorée, et le suivant est vers $1$~GHz.

\paragraph{4. Espacement moyen des raies.}
Avec les $981$ configurations accessibles, chacune portant $\simeq3.5$ niveaux sous $380$~MHz :
\[
N_{\text{raies}} \simeq 981\times3.5 \simeq 3.4\times10^3,
\qquad
\delta \simeq \frac{380~\text{MHz}}{3.4\times10^3} \simeq 110~\text{kHz}.
\]
(Sur $0$--$266$~MHz : $\simeq 3.2\times10^3$ raies, $\delta\simeq 85$~kHz.)

\paragraph{Hypothèse à garder en tête.} On suppose que les $981$ canaux convergent tous vers la même asymptote.
Ceux qui convergent vers les autres asymptotes hyperfines de $^3P_1$ (séparées de $\sim$GHz) placent leurs niveaux ailleurs :
le chiffre ci-dessus est donc un espacement \emph{minimal} (densité maximale).

\paragraph{Conclusion.} L'espacement moyen ($\sim100$~kHz) est quelques fois plus grand que la largeur naturelle ($15$~kHz),
mais bien plus petit que les largeurs de plusieurs MHz observées : élargies, ces raies se chevauchent.
\end{document}
